% !PRINTED 562
\weberpage{562}
Betrachten wir ausser dem Raume $R_\xi$ noch einen zweiten Raum $R_x$,
in dem die durch (3) oder (6) bestimmten Variablen $x$ die Coordinaten
sind, so entspricht jedem Punkte des einen Raumes ein und nur ein
Punkt des anderen Raumes. Einem endliche % [sic]
Gebiete $G_\xi$ entspricht ein endliches Gebiet $G_x$.

Um die Beziehung zwischen dem Volumen $V_\xi$ und $V_x$ zu ermitteln,
müssen wir das Integral (8) nach den Regeln über die Umformung
mehrfacher Integrale behandeln, wofür die allgemeine Formel gilt:

\weq{11}{\int\!\int\!\ldots\!\int d\xi_1\, d\xi_2 \ldots d\xi_n
= \int\!\int\!\cdots\!\int
\sqrt{\left(\Sigma\pm\frac{d\xi_1}{dx_1}\frac{d\xi_2}{dx_2}\cdots\frac{d\xi_n}{dx_n}\right)^{2}}\,
dx_1\, dx_2 \cdots dx_n\ \footnotemark.}

\footnotetext{Die allgemeine Formulirung dieser Umformung rührt von
\gesperrt{Jacobi} her (De determinantibus functionalibus; \gesperrt{Crelle},
Bd. 22, 1891, % [sic: Weber prints 1891.]
gesammelte Werke, Bd. 3). Die Ableitung findet sich in
den meisten ausführlicheren Lehrbüchern der Integralrechnung, z. B.
\gesperrt{Serret-Harnack}, Bd. 2; \gesperrt{Lipschitz}, Lehrbuch der
Analysis Bd. 2; auch \gesperrt{Baltzer}, Determinanten.}

Setzen wir darin für die Functionaldeterminante den Werth (4) oder
(5) ein, so ergiebt sich

\weq{12}{V_\xi = \sqrt{D}\ V_x,}

worin nun das Integral

\weq{13}{V_x = \int\!\int\!\cdots\!\int d x_1\, d x_2 \cdots d x_n}

sich auf das Gebiet $G_x$ bezieht.

Das Integral $V_x$ können wir als Grenzwerth einer Summe von
Elementen $\Delta x_1 \Delta x_2 \ldots \Delta x_n$ auffassen, und
wenn wir den Incrementen $\Delta x_i$ allen die Grösse $1 : t^{\frac1n}$
geben, so wird jedes dieser Elemente den Werth $1 : t$ haben.
Bedeutet $Z_t$ die Anzahl dieser Elemente, so ist $V_x$ der
Grenzwerth, dem sich das Verhältniss $Z_t : t$ bei unendlich
wachsendem $t$ nähert, also

\weq{14}{V_x = \mathrm{Lim}_{t\,=\,\infty} \frac{Z_t}{t}.}

Nun ist die Anzahl der Elemente $Z_t$ ebenso gross, wie die Anzahl
der Punkte im Inneren von $G_x$, deren Coordinaten den Ausdruck

\[
x_1 = h_1 t^{-\frac1n},\ x_2 = h_2 t^{-\frac1n},\ \ldots,\ x_n = h_n t^{-\frac1n}
\]

mit ganzzahligen $h_1, h_2, \ldots, h_n$ haben, oder auch gleich der An-
%CONT>>
